\begin{table}

\caption{\label{tab:micro-hbs-aggregation-bias-2022}Aggregation bias from harmonised HBS microdata in 2022}
\centering
\begin{tabular}[t]{llrrrrrr}
\toprule
Country & Source & Level 1 gap & Level 2 gap & Level 3 gap & L1--L2 bias & L2--L3 bias & L1--L3 bias\\
\midrule
DE & Micro HBS & 0.878 & -0.325 & NA & 1.203 & NA & NA\\
ES & Micro HBS & -0.012 & -0.107 & NA & 0.095 & NA & NA\\
FR & Micro HBS & 0.176 & -0.409 & NA & 0.585 & NA & NA\\
IT & RAS aggregated fallback & -0.269 & 1.373 & NA & -1.643 & NA & NA\\
BE & Micro HBS & -0.035 & -0.791 & -0.561 & 0.757 & -0.23 & 0.526\\
\addlinespace
NL & Micro HBS & 2.284 & 0.231 & -0.409 & 2.053 & 0.64 & 2.693\\
\bottomrule
\end{tabular}
\end{table}
\begin{flushleft}
\footnotesize{\emph{Notes:} Income quintiles are country-specific and constructed using the HBS survey weights and net household income divided by the OECD equivalence scale. Expenditure baskets are estimated directly from the 2020 Eurostat HBS scientific-use household records. Imputed rent (COICOP 042) is excluded. Inflation is measured from January to December 2022. Levels 1 and 2 use micro HBS data for DE, ES and FR; levels 1--3 use micro HBS data for BE and NL. Italy uses the existing aggregated RAS income-basket pipeline because EUR_HH099 is entirely missing in its HBS scientific-use file. A positive bias means the coarser level produces a larger Q1--Q5 inflation gap.}
\end{flushleft}
